\chapter{Roadblocks}

When working on a project, it is natural to begin with a clear and straight
path and carefully laid out plan. From the beginning, every stage seems clear,
and every task looks achievable within the anticipated timeline. However, as we
progress on the task at hand, expectations are often not met. Unforeseen
complications, external factors can emerge, challenging initial assumptions and
demanding adjustments along the way.

Throughout the course of this project, we encountered some of these moments
where things did not proceed as smoothly as originally planned. In some cases,
technical limitations, resource constraints, or changing circumstances forced
us to reevaluate our strategy. In others, new insights gained during
implementation led us to recognize better alternatives that had not been
considered initially.

In this chapter, we will explore the various obstacles we encountered, the
reasoning behind the adjustments we made, and how each decision shaped the
final result. By reflecting on these experiences, we aim to provide a
transparent account of the project’s evolution, highlighting not only the
successes but also the valuable lessons learned from overcoming difficulties
along the way.

\section{NEF Emulator}

During this phase of the project, one of the key components we relied on was
the NEF Emulator. It essentially served as a stand-in for the core’s NEF,
functioning in many ways as an actual NEF. This tool played a vital role,
letting us integrate the solution with IT's 5G core network. However, when we
first got access to it, its features were quite limited, especially for what we
needed to accomplish. Its functionalities fell short of supporting our
project’s objectives. This turned into a major technical hurdle and forced us
to invest a great deal of time into expanding its capabilities to suit our
needs.

One of the first big problems were the response and request schemas it used.
Many of them were either outdated or didn’t match the standards we were
required to follow. This became a serious issue because these schemas dictate
how systems communicate and share data. Since our work needed to comply
strictly with the accepted standards, we had no choice but to ensure every
request and response aligned with the most current specifications. Standards
like those maintained by 3GPP and CAMARA set very detailed rules for the
formats and structures that systems must use. Because of this, we had to spend
significant time analyzing the gaps between the emulator’s schemas and the
official standards, then carefully updating and correcting them to achieve full
compliance, another tough challenge involved the emulator's simulated movement
features. These were supposed to let us model the users' movement. But often,
they didn’t work properly. The simulations failed frequently, delivered
inconsistent results, or behaved unpredictably. This made it unreliable for
mobility testing, which was a crucial part of our project because we needed to
verify how APIs and the network responded to users moving around. To fix this,
we dug deep to identify the root causes of these problems. In many cases, we
had to completely rework the movement simulation functions. The new version
delivered much more reliable and stable simulations, giving us confidence in
the accuracy of our mobility scenario tests.

The biggest and most time-consuming issue was the incomplete API
implementations in the NEF Emulator. Many APIs existed only as placeholders
with no real functionality, these were for example the \emph{AnalyticsExposure}
endpoint. This posed a serious problem because these APIs were essential for
communication between the gateway and NEF. To address this, we had to develop
these APIs ourselves. We made sure they followed both the CAMARA specifications
on the gateway side and the more detailed 3GPP specifications on the NEF side.
This required a deep study of both sets of standards, careful mapping of API
behaviors, and the development of solid backend logic to handle requests and
responses as the standards required. This part of the work demanded a lot of
skill and attention to detail but ultimately resulted in a much more complete
NEF Emulator that supported all the scenarios we needed.

\section{5G Core}

The main goal of the project focused on building a smooth and dependable link
between the gateway and IT's 5G Core. The 5G Core wasn’t just another piece of
equipment. It sat right at the center of the project’s entire setup, shaping
how things were designed and built. Because so much relied on it, having the
core ready and steady played a big role in how the project moved forward. But
this dependence also brought several problems and slowdowns that had to be
handled as the work continued.

One of the biggest and most ongoing problems was simply getting access to the
5G Core in the first place. For nearly half of the project’s full timeline, the
team worked without being able to connect to the core at all. This made things
uncertain because the gateway’s features had to be developed and tested on
their own without knowing if they would work properly when finally linked to
the core. Since there wasn’t a fully connected testing setup during this early
stage, the team had to build things bit by bit while hoping it would work on
the core.

When the team finally got access to the core, more issues showed up. IT’s 5G
Core served as a shared platform, supporting many different projects from
various teams. Because of this, access to the system was controlled and limited
to avoid any problems for the other teams working on their own projects. This
setup meant a lot of back-and-forth planning with others to schedule time for
testing and integration. These coordination efforts often took careful planning
because the core was frequently booked for other work. As a result, more delays
followed before the solution could be fully tested and connected with the rest
of the system.

Another problem came from the way the system was built. The NEF was supposed to
interact with a Slice Manager API, which acted as a go-between, connecting with
Huawei’s 5G Core. But during the project, the IT staff member who created and
ran the Slice Manager API left IT. This sudden exit left a big knowledge gap.
The remaining documentation wasn’t much help, as it was incomplete and
sometimes even wrong, mentioning fields and parameters that either didn’t exist
or worked differently than described. Because of this, the team had to spend a
lot of time figuring things out on their own to make the API work properly for
integration. On top of that, documentation for the core itself arrived later
than expected and often lacked enough detail to solve problems quickly, making
the process even harder.

Besides all the technical issues, the project ran into a completely unexpected
outside problem, a countrywide blackout. The blackout only lasted a day, but
its effects dragged on much longer. Power loss hit the normal business
operations and directly impacted the 5G Core’s systems. Several of its parts
needed to be checked, fixed, and restarted to safely bring them back online.
This recovery effort led to more delays and added more time when the team
couldn’t access the core for integration, testing, and validation work.
